\emph{(i)} A tower of $(1,1)$ ascents is the sequence $(p+k, q+k)$, $k \geq 0$, over a ground with $\min(p,q) = 0$, that is a definite $(m,0)$ or $(0,m)$. Then $\min(p+k,q+k) = k$, which equals $1$ at $k = 1$ only, so the unique causal floor is $(m+1,1)$ or $(1,m+1)$, of dimension $m+2$. The residue $p - q$ is unchanged by the ascent, and $\Cl(p+k,q+k) \cong \Mat(2^{k}, \Cl(p,q))$ (Fact \claimref{dark-matter.imported-eight-towers-mod-8-clock}), so residue, Morita class and division-ring type are constant along each tower. The ground $(m,0)$ has residue $m \bmod 8$ and $(0,m)$ has residue $-m \bmod 8$; so the causal signatures of residue $r$ are the causal floors of the towers with $m \equiv r$, respectively $m \equiv -r$, modulo $8$, the least $m$ being $\min(r, 8-r)$ (with $0$ for $r = 0$). This reproduces the census of Proposition \claimref{dark-matter.towers-are-residues-single-clock-filter}(i), and adds that the causal filter picks exactly one floor in each tower. The check enumerates every ground of dimension at most twelve.

\emph{(ii)} The objects are those of Definition \claimref{obj.dark-tower-matter.ideals-fields-momentum}, and the momentum is that of Premise \claimref{dark-tower-matter.h1-momentum}, read at dimensions two and six as at three and five; this extension is a hypothesis of the statement. At $(1,1)$ the even algebra is spanned by $1$ and $b = e_{t}e_{a}$ with $b^{2} = -Q(e_{t})Q(e_{a}) = 1$, so $\Cl^{0}(V) = \R f \oplus \R(1-f)$ with $f = \tfrac12(1+b)$, and $S = \Cl^{0}(V)f = \R f$. For $\psi = \lambda f$, $j(\psi) = \lambda^{2} j(f)$. Since $e_{t}$ anticommutes with $b$, $e_{t}f = (1-f)e_{t}$ and $\rev{f} = 1 - f$, so $f e_{t}\rev{f} = f e_{t}(1-f) = f e_{t}$ and $j(f) = \tfrac12(e_{t} - Q(e_{t})e_{a})$, a null vector. Every total momentum $\sum_{i}\lambda_{i}^{2} j(f)$ is a nonnegative multiple of $j(f)$, hence null, and $m^{2} = 0$ for every finite set of constituents from $S$. The same computation for $1 - f$ gives $\tfrac12(e_{t} + Q(e_{t})e_{a})$, the other null ray, and the sum of the two has $m^{2} = Q(e_{t})Q(e_{t}) = 1$. The two ideals are not isomorphic as $\Cl^{0}(V)$-modules, $f$ and $1 - f$ being distinct central idempotents. The definition's $f = \tfrac12(1 + e_{t}e_{a})$ selects one of them by the orientation of $e_{a}$ (replacing $e_{a}$ by $-e_{a}$ exchanges $f$ and $1 - f$), so the single spinor ideal is a chirality choice; the conclusion of (ii) is relative to it.

\emph{(iii)} At $(1,5)$ and $(5,1)$ the pseudoscalar $I$ of $V$ is even, central in $\Cl^{0}(V)$ and satisfies $I^{2} = 1$; it commutes with $b$ and is invariant under the adjoint, so $f = \tfrac14(1+b)(1 \pm I)$ is an adjoint-invariant idempotent. The check finds $\dim_{\R} S = 8$ and $\dim_{\R} f\Cl^{0}(V)f = 4$, and it exhibits two anticommuting elements of the corner with square a negative multiple of $f$, so the corner is $\HH$, a division algebra, $f$ is primitive, and the field is $\HH$, the division algebra of residue $5$ as in Theorem \claimref{dark-tower-matter.constituents-field}. Writing $\psi = \sum_{i} x_{i}s_{i}$ on a basis $s_{1}, \ldots, s_{8}$ of $S$, $j(\psi)$ is a quadratic form in $x$ with vector coefficients $J_{il}$, and $Q(j(\psi)) = \sum Q(J_{il}, J_{km})\,x_{i}x_{l}x_{k}x_{m}$; the check collects the coefficient of every quartic monomial and finds each zero, so $Q(j(\psi)) = 0$ as a polynomial identity over $\Q$, not on samples. Two basis constituents have total momentum with $m^{2} = \tfrac14 \neq 0$. Hence mass begins at two constituents. The same routine reproduces Theorem \claimref{dark-tower-matter.where-mass-begins} at $(1,3)$ and $(2,1)$ (single momenta null, pairs massive) and at $(1,4)$ (single momenta not all null), and the fields of Theorem \claimref{dark-tower-matter.constituents-field} at those signatures.

\emph{(iv)} Proposition \claimref{dark-matter.unit-ball-rule} gives each populated dark tower the share $V_{3}/V_{D}$ with $D$ its space dimension; $V_{3} = \tfrac{4\pi}{3}$ and $V_{5} = \tfrac{8\pi^{2}}{15}$, so $V_{3}/V_{5} = \tfrac{5}{2\pi}$, and $1 + \tfrac83 + \tfrac{16}{3\pi} + \tfrac{5}{2\pi} = \tfrac{11}{3} + \tfrac{47}{6\pi}$. With $\tfrac{333}{106} < \pi < \tfrac{355}{113}$ the check shows that $\tfrac{11}{3} + \tfrac{16}{3\pi}$ lies inside $5.364 \pm 0.065$ and that $6.160 < \tfrac{11}{3} + \tfrac{47}{6\pi} < 6.161$, between $12$ and $13$ standard errors above $5.364$. Proposition \claimref{dark-matter.status-rule} records that the unit-ball relation was found by solving for the observed value and that its clauses were chosen with it known, so the agreement is a fit; the emptiness of residue $4$ is among those clauses. Which towers are populated is observation-level (Observation \claimref{dark-matter.which-towers-are-populated}); the statement is only that residue $4$'s emptiness, on which the match depends, is not a consequence of the content in (iii). The lowest-causal choice itself is the minimality reading that Premise \claimref{dark-matter.p3-nearest-dimension-representatives} cites as its support, applied inside the causal filter; it is not derived here.
